\documentclass{article}

\begin{document}

Basically, here goes everthing i need to make sure 
i talk about in my thesis

\begin{itemize}
\item Revisiter avec un oeil critique les r\'esultats obtenus dans 
      le papier de 2016
\item Se pencher sur la fonction coulombienne mais aussi le paquet 
      d'onde ionis\'e pour d\'eceler une eventuelle erreur dans la 
      gestion des projections 
\item Comparer CNCL, RKLA, et CNFA du point de vue qualit\'e 
      et du point de vue performance 
\item \'Etudier l'incidence du nombre de sturmiennes sur la convergence 
      si possible pour tous les sch\'emas, sur les spectres normaux et
      semi-logarithmiques pour tous les auteurs
\item \'Etudier l'incidence du nombre de moments angulaires sur la convergence
       si possible pour tous les sch\'emas, sur les spectres normaux et
       semi-logarithmiques
\item Comparer les r\'esultats pour les sch\'emas de propagation de type 
      Crank-Nicholson pour ivnufa=0 et 1
\item Revoir le sch\'ema de Nuhruda-Faisal et ma\^itriser les conditions 
      d'applicabilit\'e de leur algorithme
\item Quelle est la forme de la matrice d'interaction dans la base 
      atomique? Si elle est pleine, alors essayer de 
      tirer partie de la structure a bandes de cette matrice 
      pour la propagation.
\item Quelle est la normalisation de la paquet d'onde ionis\'e que nous
      utilisons? Ce paquet d'onde est il normalis\'e sur les \'energies 
      ou sur les impulsions?
\item Le param\`etre d'asym\'etrie beta...
\item Pour des r\'esultats qui convergent et dont les spectres 
      co\"incident, mesurer la d\'eviation sur la distribution 
      de probabilit\'e, le paquet d'onde (parties r\'eelle et imaginaire)
\item La troncature du potentiel pour une comparaison avec le 
      mod\`ele des potentiels s\'eparables. (tr\`es important)
\item Obtenir quelques r\'esultats avec la quadrature de Gauss-Laguerre 
      pour le traitement du potentiel
\item G\'en\'erer des r\'esultats dans la base des B-splines dans 
      son actuelle implementation.
\item Faire varier le nombre Z pour modeliser les hydrog\'eno\"ides.
\item Faire l'\'etude sur le nombre de nsteps de steps pour les moteurs 
      de type Crank Nicholson.
\item La discusion d'Hamido (p.73 de sa th\`ese) autour des 
      oscillations des distributions 
      quand on extrait pas les \'etats li\'es avant de faire la 
      projection pour obtenir les distributions. \`A ce sujet, 
      utiliser la projection sur la fonction coulombienne entrante, 
      non pas du paquet d'onde final, mais du paquet d'onde "ionis\'e", 
      c'est \`a dire, dont les \'etats li\'es ont \'et\'e extraits.   
\end{itemize}
\end{document}
